% Generated draft: requires mathematical and source review.
\documentclass[11pt]{article}
\usepackage[margin=25mm]{geometry}
\usepackage{fontspec}
\setmainfont{Latin Modern Roman}
\usepackage{amsmath,amssymb,mathtools,unicode-math}
\setmathfont{Latin Modern Math}
\usepackage{enumitem,needspace,xurl,hyperref}
\setlength{\parindent}{0pt}
\setlength{\parskip}{5pt}
\setlength{\emergencystretch}{3em}
\newcommand{\N}{\mathbb{N}}
\newcommand{\Z}{\mathbb{Z}}
\newcommand{\Q}{\mathbb{Q}}
\newcommand{\R}{\mathbb{R}}
\newcommand{\C}{\mathbb{C}}
\newcommand{\F}{\mathbb{F}}
\newcommand{\A}{\mathbb{A}}
\newcommand{\PP}{\mathbb{P}}
\newcommand{\E}{\mathbb{E}}
\newcommand{\eps}{\varepsilon}
\newcommand{\dd}{\,\mathrm{d}}
\begin{document}
\section*{UnsolvedMath ID: 10400158}
\textbf{Source problem code:} catalogue entry
\subsection{Definitions and mathematical statement}
For the primes \(p\) and genera \(g\) in the source's cyclotomic TQFT setup, let \(\zeta_p\) be a primitive \(p\)-th root of unity and \(R_p=\Z[\zeta_p]\). Write \(\Sigma_g\) for the closed oriented genus-\(g\) surface and retain the source's notation \(V_p\) and \(V_p^{[j]}(M)\), whose constructions are not included. The displayed exponent below requires the intended odd-prime setting; the exact lower bound on \(p\) is omitted.

\emph{Problem (T. Kerler, cyclotomic integer TQFTs).}
\begin{enumerate}
\item Find explicit, computable bases exhibiting \(V_p(\Sigma_g)\) as free \(R_p\)-modules.
\item Show that \(V_p\) extends to all cobordisms as a half-projective TQFT with parameter
\[
 x=(\zeta_p-1)^{(p-3)/2}\in R_p.
\]
The original domain, extra structures on cobordisms, and half-projective composition rule are not specified in the excerpt.
\item Determine the structure of \(V_p^{[j]}(M)\) and the extent to which these objects lift from \(\F_p\) to \(\Z\), analogously to Ohtsuki invariants of closed \(3\)-manifolds. The source expressly requires a lift to depend on \(p\), since the relevant vector-space dimensions do. It also requires further quotients: the irreducible TQFTs over \(\Z\) do not have the required dimensions, while their reductions to \(\F_p\) become reducible.
\item Find a universal TQFT combining all \(V_p\), at least perturbatively, into one theory. The source explicitly calls this request vague and suggests an infinite filtered space with finite graded components; it gives no universal property or compatibility maps.
\end{enumerate}
The source reports that items (1)--(3) were mostly carried out for \(p=5\) in [213], that little was known for primes \(p\geq7\), and that Gilmer found some explicit bases for \(g=1\), while the situation for \(g\geq3\) was unknown. It supplies no corresponding assertion for \(g=2\). It states that (2) would imply that the quantum order of [92] is an upper bound for the cut-number of a \(3\)-manifold, and that a closely related statement would give a different proof that Ohtsuki invariants are of finite type. Neither the quantum order nor that related statement is defined here.

The trailing context defines a TQFT \(V:\operatorname{Cob}\to R\text{-}\operatorname{mod}\) to be homological when each induced mapping-class-group representation
\[
 V[g]:\Gamma_g\longrightarrow\operatorname{GL}_R(V(\Sigma_g))
\]
factors through the homology-action quotient \(\Gamma_g\twoheadrightarrow\operatorname{Sp}(2g,\Z)\). It calls such a theory strictly homological when each resulting representation is ``algebraic, i.e. either faithful or zero.'' This quoted identification is not a standard definition of an algebraic representation and is retained as an unresolved source qualification, rather than used as an equivalence. The source gives the Lefschetz components \(V^{(j)}\) of the Frohman--Nicas TQFT over \(\Z\) as examples [128, 214]. It forms \(\mathcal Q_0\) by taking direct sums of these components, says that all TQFTs in [110] belong to \(\mathcal Q_0\), and forms a larger family \(\mathcal Q_*\) using tensor products and their irreducible summands. These are contextual claims, not additional numbered construction requirements.
\subsection{Short English statement}
Kerler asks for computable integral bases, an extension to all cobordisms with the specified half-projective parameter, integral lifts of the mod-\(p\) pieces, and a perturbative theory unifying the prime-dependent TQFTs. The source records substantial \(p=5\) progress, limited genus-one bases, and several unresolved definitions.
\subsection{Sources}
\begin{itemize}
\item Local UnsolvedMath registry, permanent ID 10400158, source code catalogue entry; attribution T. Kerler, ``Cyclotomic integer TQFT's.'' Supplied title is truncated after ``(T.''.
\item Ohtsuki, \emph{Problems on invariants of knots and 3-manifolds} (2002), Problem 8.12, PDF page 126. Supplied URL: \url{https://msp.org/gtm/2002/04/gtm-2002-04-024s.pdf}; recorded access 2026-07-29; extraction pdf-font-structure. The collection distinguishes 231 numbered bold problem, conjecture, and question headings.
\item Source references [213] for the \(p=5\) program, [92] for quantum order, [128, 214] for Frohman--Nicas Lefschetz components, and [110] for examples in \(\mathcal Q_0\); their bibliographic entries are not supplied. Gilmer is credited with some explicit genus-one bases.
\item The supplied registry labels the item partially\_solved and PARTIAL-PROGRESS; the background describes its original presentation as an open problem with current status NEEDS\_REVIEW. The research summary says higher primes and higher genera remain open; these are supplied assessments, not an independently verified current audit. Redistribution terms for the publicly accessible author/publisher copy require review.
\end{itemize}
\end{document}
